\documentclass[12pt,a4paper]{article}

\usepackage[margin=2cm]{geometry}
\usepackage{amsmath}
\usepackage{enumitem}

\setlength{\parindent}{0pt}
\setlength{\parskip}{6pt}

\begin{document}

\begin{center}
    {\Large\textbf{Statics -- Resolving Forces A}}\\[4pt]
    \textbf{A-Level Physics}
\end{center}

For each question, draw a force diagram and resolve the forces horizontally and vertically.

Where required, take \(g = 9.81\,\mathrm{m\,s^{-2}}\).

\vspace{0.4cm}

\begin{enumerate}[leftmargin=*, itemsep=1em]

\item
A \(20\,\mathrm{kg}\) box rests on a horizontal floor. A rope pulls the box with a tension of \(75\,\mathrm{N}\) at \(30^\circ\) above the horizontal.

Calculate the normal reaction force exerted by the floor.

\vspace{0.8cm}

\item
A \(35\,\mathrm{kg}\) crate rests on a horizontal floor. A force of \(120\,\mathrm{N}\) acts at \(25^\circ\) below the horizontal.

Calculate the normal reaction force exerted by the floor.

\vspace{0.8cm}

\item
A box remains stationary on a horizontal floor. It is pulled by a rope with a tension of \(140\,\mathrm{N}\) acting at \(40^\circ\) above the horizontal.

Calculate the frictional force acting on the box.

\vspace{0.8cm}

\item
A \(22\,\mathrm{kg}\) box rests on a horizontal floor. A rope pulls it with a tension of \(90\,\mathrm{N}\) at \(35^\circ\) above the horizontal. The box remains stationary.

Calculate:

\begin{enumerate}[label=(\alph*)]
    \item the normal reaction force from the floor;
    \item the frictional force acting on the box.
\end{enumerate}

\vspace{0.8cm}

\item
A \(40\,\mathrm{kg}\) crate rests on a horizontal floor. A force of \(160\,\mathrm{N}\) is applied at \(20^\circ\) below the horizontal. The crate remains stationary.

Calculate:

\begin{enumerate}[label=(\alph*)]
    \item the normal reaction force;
    \item the frictional force.
\end{enumerate}

\vspace{0.8cm}

\newpage

\item
A \(12\,\mathrm{kg}\) object is held stationary against a vertical wall by a force of \(200\,\mathrm{N}\) acting at \(40^\circ\) above the horizontal towards the wall.

Calculate:

\begin{enumerate}[label=(\alph*)]
    \item the normal reaction force exerted by the wall;
    \item the frictional force acting between the object and the wall.
\end{enumerate}

\vspace{0.8cm}

\item
A \(25\,\mathrm{kg}\) box rests on a horizontal floor. A force of \(180\,\mathrm{N}\) acts at \(50^\circ\) below the horizontal to the right. A second force acts horizontally to the left and keeps the box in equilibrium.

Calculate:

\begin{enumerate}[label=(\alph*)]
    \item the horizontal force acting to the left;
    \item the normal reaction force from the floor.
\end{enumerate}

\vspace{0.8cm}

\item
A \(30\,\mathrm{kg}\) crate rests on a horizontal floor. It is pulled by a cable with a tension of \(150\,\mathrm{N}\) acting at \(35^\circ\) above the horizontal to the right. A horizontal force acts to the left and keeps the crate stationary.

Calculate:

\begin{enumerate}[label=(\alph*)]
    \item the horizontal force acting to the left;
    \item the normal reaction force from the floor.
\end{enumerate}

\vspace{0.8cm}

\item
A \(24\,\mathrm{kg}\) box rests on a horizontal floor. A rope pulls the box with a tension of \(100\,\mathrm{N}\). The box remains stationary and the normal reaction force from the floor is \(190\,\mathrm{N}\).

Calculate:

\begin{enumerate}[label=(\alph*)]
    \item the vertical component of the tension;
    \item the angle that the rope makes above the horizontal.
\end{enumerate}

\vspace{0.8cm}

\item
A \(30\,\mathrm{kg}\) crate rests on a horizontal floor. A force of \(140\,\mathrm{N}\) acts at an unknown angle below the horizontal. The normal reaction force from the floor is \(370\,\mathrm{N}\).

Calculate:

\begin{enumerate}[label=(\alph*)]
    \item the vertical component of the applied force;
    \item the angle that the force makes below the horizontal.
\end{enumerate}

\end{enumerate}

\newpage

\begin{center}
    {\Large\textbf{Statics -- Answer Key}}
\end{center}

\vspace{0.5cm}

\begin{enumerate}[leftmargin=*, itemsep=1em]

\item
Normal reaction force: \(158.7\,\mathrm{N}\)

\item
Normal reaction force: \(394.0\,\mathrm{N}\)

\item
Frictional force: \(107.2\,\mathrm{N}\)

\item
Normal reaction force: \(164.2\,\mathrm{N}\) \\
Frictional force: \(73.7\,\mathrm{N}\)

\item
Normal reaction force: \(447.1\,\mathrm{N}\) \\
Frictional force: \(150.4\,\mathrm{N}\)

\item
Normal reaction force: \(153.2\,\mathrm{N}\) \\
Frictional force: \(10.8\,\mathrm{N}\) downward

\item
Horizontal force to the left: \(115.7\,\mathrm{N}\) \\
Normal reaction force: \(383.1\,\mathrm{N}\)

\item
Horizontal force to the left: \(122.9\,\mathrm{N}\) \\
Normal reaction force: \(208.3\,\mathrm{N}\)

\item
Vertical component of tension: \(45.4\,\mathrm{N}\) \\
Angle above the horizontal: \(27.0^\circ\)

\item
Vertical component of applied force: \(75.7\,\mathrm{N}\) \\
Angle below the horizontal: \(32.7^\circ\)

\end{enumerate}

\end{document}